% GENERATED, edit the French unit or the translation table
% GENERATED from content/data/enrichment.json by mxbuild.enrichment; do not edit.
\mxunit{enrichment-modular-notes}{Congruences and last digits}
\mxsubtitle{High school · number theory \textperiodcentered{} Lesson and method \textperiodcentered{} 03}

\begin{prerequisites}
Divisibility, powers and Euclidean division.
\end{prerequisites}

\begin{keypoint}[Ideas and method]
$a \equiv  b$ (mod m), with $m \geq  2$, means that m divides $a - b$. They have the same remainder upon division by m. For example, $-1 \equiv  6$ (mod 7).

Congruences are preserved by addition and multiplication, hence by positive integer powers. Work modulo 10 to find a last digit.

A factor cannot always be cancelled: $2 \times  1 \equiv  2 \times  4$ (mod 6), but 1 is not congruent to 4 modulo 6. A factor coprime to the modulus can be cancelled.
\end{keypoint}

\begin{example}
The last digits of $3^{1}, 3^{2}, 3^{3}, 3^{4}$ are $3,9,7,1$. The cycle repeats because $3^{4} \equiv  1$ (mod 10). Since $2026 = 4 \times  506 + 2$, the last digit of $3^{2026}$ is 9.
\end{example}

\begin{activity}[Explore]
Build tables of squares modulo 3, 4 and 5. Identify impossible remainders.
\end{activity}

